\chapter{相关技术与研究现状}

\section{布控球与 PTZ 云台控制技术}

PTZ 是 Pan、Tilt、Zoom 的缩写，分别表示水平旋转、俯仰旋转和变焦控制。布控球通常把摄像机、云台、通信模块和电源模块集成在一台便携式设备中，通过有线或无线网络向上位机传输视频流，并接收上位机下发的 PTZ 控制指令。与固定摄像机相比，PTZ 设备能够通过主动转动扩大监控范围，也能在检测到目标后调整姿态，使目标尽量保持在画面中心附近。对于移动式监控设备来说，PTZ 控制不仅决定画面覆盖范围，也会直接影响目标跟踪质量和后续视频分析效果。

传统 PTZ 控制方法大致可以分为人工遥控、预置点巡航和基于规则的自动控制。人工遥控依赖操作人员经验，比较适合低频操作场景，但在目标快速移动时响应速度往往不足。预置点巡航能够覆盖多个固定区域，却很难对随机运动目标做出灵活反馈。基于视觉反馈的 PTZ 控制会根据目标位置估计和控制策略自动生成云台指令，因此更适合动态目标跟踪场景。Petrov 针对视觉目标跟踪提出了自适应云台控制方法，说明目标位置误差与云台运动控制之间可以建立反馈关系\cite{petrov_adaptive_ptz}。Huynh 等研究了基于视觉反馈的反步控制方法，强调了控制模型对跟踪稳定性的影响\cite{huynh_vision_backstepping_2021}。Guo 等针对小型惯性稳定云台开展视觉检测与跟踪研究，为复杂运动场景下的稳定控制提供了参考\cite{guo_vision_based_2021}。

从控制系统角度看，视觉 PTZ 跟踪的关键，是把图像平面上的目标偏差转换为设备空间中的云台动作。理想情况下，系统可以通过摄像机标定、焦距估计和云台角速度模型建立较精确的映射关系；但在实际布控球设备上，SDK 接口往往只提供离散方向控制和速度参数，并且机械响应本身也会存在一定迟滞。因此，工程系统更常采用阈值、脉冲、冷却和平滑等可解释策略，以较低建模成本获得可以接受的跟踪效果。

在实际工程系统中，PTZ 控制还需要考虑设备协议、指令频率、单次转动幅度和机械响应延迟等因素。如果控制指令过于频繁，云台容易出现抖动；如果控制幅度过大，又可能越过目标中心并造成反复修正。因此，本文采用偏差阈值、比例脉冲和平滑滤波相结合的策略，在不建立复杂物理模型的前提下实现可调试、可部署的闭环控制。该策略与真实设备调试过程比较匹配，能够通过调整少量参数逐步平衡响应速度与稳定性。

对于本文所研究的布控球场景而言，控制策略需要在精确性和可实现性之间取舍。若建立完整的相机标定模型和云台动力学模型，理论上可以得到更准确的控制量，但实际设备的 SDK 接口、安装姿态和机械响应并不总是完全透明，参数标定成本也较高。相反，基于图像偏差的规则控制虽然精度有限，但参数含义清楚，便于在真实设备上快速调试。本文选择后一种方式，主要是为了在原型阶段优先验证完整闭环是否可行。

\section{轻量级目标检测技术}

目标检测是智能视频监控系统中的核心感知能力，主要任务是在图像中定位目标并判断类别。早期检测算法通常依赖手工特征和滑动窗口，计算量较大，泛化能力也比较有限。深度学习兴起以后，两阶段检测方法和单阶段检测方法逐渐成为主流方向。Faster R-CNN 通过区域建议网络提高了候选区域生成效率，使检测精度得到明显提升\cite{ren_faster_rcnn_2015}；SSD 将多尺度特征图用于单阶段检测，在速度和精度之间取得了较好的平衡\cite{liu_ssd_2016}。YOLO 系列进一步把目标定位和分类统一为单阶段预测任务，具有结构简洁、推理速度快和部署生态成熟等特点\cite{redmon_yolo_2016,redmon_yolov3_2018}。在需要实时反馈控制的场景中，检测算法不能只追求较高准确率，还要保证输出延迟足够低，否则云台动作会落后于目标运动。

嵌入式布控球系统对目标检测提出了实时性和资源占用方面的双重要求。一方面，检测模块需要及时输出目标坐标，否则 PTZ 控制会滞后于目标运动；另一方面，边缘设备的算力和功耗有限，模型参数量、输入尺寸和推理精度都需要权衡。轻量化网络在移动视觉任务中已经得到较多研究，例如 MobileNets 通过深度可分离卷积减少计算量\cite{howard_mobilenets_2017}，MobileNetV2 进一步利用倒残差结构和线性瓶颈提升了移动端特征表达效率\cite{sandler_mobilenetv2_2018}。这些研究说明，在边缘设备上部署视觉模型时，模型结构、输入尺寸和推理后端需要整体考虑。YOLOv8n 属于 YOLOv8 系列中的轻量级模型，适合在资源受限设备上快速验证人体检测和跟踪闭环。Ultralytics 提供的 Python 接口降低了模型加载、推理和结果解析的工程门槛\cite{ultralytics_yolov8}。本文系统将输入尺寸设置为较小规模，并结合置信度阈值过滤人体类别，使检测模块能够在实时性和可用性之间取得折中。

模型压缩技术能够进一步改善嵌入式部署效果。Han 等提出的深度压缩方法将剪枝、量化和编码结合起来，证明神经网络可以在较小精度损失下降低存储和计算开销\cite{han_deep_compression_2016}。Batch Normalization 和残差结构等基础网络技术也在深度模型训练和表达能力提升中发挥了重要作用\cite{ioffe_batchnorm_2015,he_resnet_2016}，为轻量模型训练和迁移提供了基础。在布控球场景中，可以通过 INT8 量化、TensorRT 加速、输入尺寸自适应和检测频率动态调整等方法继续降低推理延迟和资源占用\cite{nvidia_tensorrt_2024}。同时，检测频率也不一定要始终与视频帧率完全一致，系统可以根据目标运动速度、设备算力和控制需求动态决定是否跳帧推理，从而降低边缘端负载。

\section{目标跟踪与人体重识别技术}

目标检测只能给出单帧中的目标位置，不能直接保证跨帧身份一致。在多人场景中，如果系统始终选择面积最大的检测框，跟踪目标可能会随着人群距离变化而频繁切换。目标跟踪和人体重识别技术主要用于解决跨帧关联问题。常见跟踪方法包括基于运动模型的 Kalman 滤波、基于检测关联的 SORT/DeepSORT，以及基于外观特征的人体 ReID。Kalman 滤波为状态估计和运动预测提供了经典方法基础\cite{kalman_filter_1960}，SORT 则将检测结果、Kalman 预测和数据关联结合起来，形成了较为高效的在线多目标跟踪框架\cite{bewley_sort_2016}。Wojke 等提出的 DeepSORT 将深度外观描述子引入目标关联过程，提升了遮挡和交叉场景下的身份保持能力\cite{wojke_deepsort_2017}。Luo 等提出了带批归一化颈部结构的强基线方法，说明外观特征学习对行人重识别性能具有重要影响\cite{luo_reid_baseline_2020}。

对于布控球自动跟踪系统而言，目标身份一致性会直接影响云台运动的连续性。当画面中两个行人交叉，或者人物距离摄像机的远近发生变化时，单纯依靠检测框面积选择目标，可能会导致控制目标突然切换，进而造成云台方向变化和用户观感不稳定。因此，即使不引入复杂的深度 ReID 网络，系统也需要一种成本较低的目标记忆机制，用来在短时间遮挡、出入画面和多人并存时维持基本的目标连续性。

完整的深度 ReID 模型通常需要额外的神经网络推理，对边缘实时系统而言会增加计算负担。本文以工程可用性为目标，设计了轻量级人体 ReID 记忆模块，使用人体框几何比例和区域平均颜色构成低维特征，并通过余弦相似度完成短时匹配。该方法定位于实时跟踪场景下的轻量身份辅助，能够在不明显增加算力开销的情况下，改善目标短时间离开画面、重新入镜以及多人同时出现时的目标连续性。其设计思路是优先满足原型系统的实时闭环要求，把身份记忆作为目标选择层的辅助能力，使系统复杂度与边缘端实时性保持平衡。

从系统角度看，目标跟踪并不一定要求每一帧都得到完全准确的身份判断，更重要的是在控制过程中避免目标频繁跳变。对于布控球自动跟踪而言，短时间内保持同一目标的连续性，往往比追求复杂场景下的长期身份识别更加符合原型系统需求。基于这一考虑，本文没有直接引入较重的 ReID 网络，而是采用轻量特征作为过渡方案，在保证实时性的前提下增强多目标场景中的目标选择稳定性。

\section{边缘智能与嵌入式部署}

边缘智能强调将计算任务下沉到靠近数据源的终端或边缘节点，从而减少数据上传量和云端处理压力。对于视频监控场景，本地推理可以降低延迟、节省带宽，也能减少敏感视频持续上传带来的隐私风险。NVIDIA Jetson、Intel Movidius 等嵌入式 AI 平台的发展，使轻量级深度学习模型在移动终端上的部署成为可能。Dagli 等关于异构 SoC 上 DNN 并发执行的研究表明，边缘设备的内存竞争、加速器调度和端到端延迟均会影响 DNN 应用性能\cite{dagli_shared_memory_2024}。这也说明，边缘智能系统并不只是模型部署问题，还涉及视频解码、内存传输、线程调度和外设控制等完整链路。

本文当前原型阶段主要关注算法、设备接口与控制闭环的可行性验证，也具备面向 Jetson 等嵌入式平台开展部署优化的基础。迁移过程中需要重点处理几个方面：一是设备 SDK 与边缘计算平台之间的接口兼容性，二是 YOLO 模型推理后端对 GPU、TensorRT 或 CPU 的适配，三是视频解码、检测线程和 PTZ 控制线程之间的实时调度。此外，嵌入式部署还需要考虑供电、散热、开机自启动、异常恢复和远程维护等工程因素，这些因素都会影响系统从实验原型走向长期运行设备的可用性。

\section{现有研究不足与本文工作定位}

总体来看，现有研究在目标检测、PTZ 控制和目标跟踪方面已经积累了较多方法，为布控球智能跟踪系统提供了较好的技术基础。面向真实布控球设备进行工程化集成时，需要把这些方法放到统一链路中综合考虑。一是实际设备 SDK 调用和码流解码过程需要与算法输入格式衔接；二是轻量模型在边缘端运行时需要兼顾帧率、精度和资源占用；三是 PTZ 控制要同时处理检测抖动、机械响应和多人目标切换；四是系统还需要提供便于调试的界面和部署模式。换言之，布控球智能跟踪并不是单一算法模块能够独立完成的问题，而是一个由感知、决策、控制和交互共同构成的实时系统问题。

本文的工作定位不是提出新的通用目标检测网络，而是面向真实布控球设备构建完整的感知--决策--控制闭环。本文重点处理 SDK 集成、实时目标检测、轻量目标记忆、平滑 PTZ 控制和可视化交互等问题，为系统性能优化和嵌入式迁移提供基础。通过将成熟模型与设备接口、控制规则和交互界面结合，本文说明了低成本边缘智能方案在布控球目标跟踪中的可行性，并形成了具有可复用价值的系统框架。
